\TieudeT{PHONG TỎA VẬT LÍ 11 - DAO ĐỘNG VÀ SÓNG}
\subsubsection*{PHẦN I. Thí sinh trả lời câu hỏi từ câu 1 đến câu 18. Mỗi câu hỏi thí sinh chỉ chọn một phương án}
\setcounter{ex}{0}
	\begin{ex} %[3P1Y1-1][Loai1] 
		Một vật dao động điều hòa trên trục Ox quanh vị trí cân bằng O. Gọi $A,\ \omega$ và $\varphi$ lần lượt là biên độ, tần số góc và pha ban đầu của dao động. Biểu thức li độ của vật theo thời gian t là
		\choice
		{\True $x = A\cos\left(\omega t + \varphi\right)$}
		{$x = \omega \cos\left(t\varphi +A\right) $}
		{$x = t\cos\left(\varphi A + \omega\right)$}
		{$x = \varphi \cos\left(A\omega +t\right)$}
		\loigiai{
			Biểu thức li độ của vật theo thời gian là $x = A\cos\left(\omega t + \varphi\right)$.
		}
		%\haidong{5}
	\end{ex}
	\begin{ex}
		Một chất điểm dao động điều hoà có tần số góc $\omega=10\pi\ rad/s$. Tần số của dao động là
		\choice
		{\True $5\ Hz$}
		{$10\ Hz$}
		{$20\ Hz$}
		{$5\pi\ Hz$}
		\loigiai{
			\begin{itemize}
				\item Tần số của dao động là: $f=\dfrac{\omega}{2\pi}=5\ Hz$.
			\end{itemize}
		}	
		%\haidong{5}
	\end{ex}
	\begin{ex}%[3P1Y2-1][Loai1]
		Trong quá trình dao động, vật có tốc độ cực đại khi vật
		\choice
		{\True đi qua vị trí cân bằng}
		{đi qua vị trí cân bằng theo chiều dương}
		{biên âm ($x = -A$)}
		{biên dương ($x=A$)}
		\loigiai{
			Trong quá trình dao động, vật có tốc độ cực đại khi vật đi qua vị trí cân bằng.
		}
		%\haidong{5}
	\end{ex}
	%%\loai{Cho khoảng cách từ vận tới vị trí cân bằng, tìm v}
	\begin{ex}%[3P1B2-2][Loai1]
		Một vật dao động điều hòa trên trục $Ox$. Khi vật qua vị trí cân bằng, tốc độ của nó là $20\ cm/s$. Khi vật ở biên, gia tốc của vật có độ lớn là $0,8\ m/s^2$. Khi vật cách vị trí cân bằng $4\ cm$ thì nó có tốc độ
		\choice
		{\True $12\ cm/s$}
		{$20\ cm/s$}
		{$25\ cm/s$}
		{$35\ cm/s$}
		\loigiai{
			\begin{itemize}
				\item Ta có: $v_{\max} = 20\ cm/s;\ a_{\max} = 0,8\ m/s^2 = 80\ cm/s^2 \Rightarrow \omega  = 4\ rad/s;\ A = 5\ cm$.
				\item Khi $\left| x \right|=4\ cm$ thì $x^2+\dfrac{v^2}{\omega^2}=A^2\Rightarrow {{4}^2}+\dfrac{v^2}{{{4}^2}}={{5}^2}\Rightarrow \left| v \right|=12\ cm/s$.
			\end{itemize}
		}
		
		%\haidong{13}
	\end{ex}
	\begin{ex}%[3P1K4-1][Loai1] 
		Phương trình của một dao động điều hoà của một chất điểm có dạng \linebreak $x = 6\sin\left(10\pi t + \dfrac{\pi}{6}\right)\ cm$. Khi pha dao động bằng $-\dfrac{\pi}{4}$ thì li độ của chất điểm là
		\choice
		{$x = \sqrt{2}\ cm$}
		{\True $x = 3\sqrt{2}\ cm$}
		{$x = 3\ cm$}
		{$x = 6\ cm$}
		\loigiai{
			\begin{itemize}
				\item Ta có: $x = 6\sin\left(10\pi t+ \dfrac{\pi}{6}\right) = 6\cos\left(10\pi t- \dfrac{\pi}{3}\right)$.
				\item Khi pha dao động bằng $-\dfrac{\pi}{4}$ thì $x = 6\cos\left(-\dfrac{\pi}{4}\right) = 3\sqrt{2}\ cm$.
			\end{itemize}
		}
		%\haidong{10}
	\end{ex}
	%%\loai{Cho khoảng thời gian ngắn nhất giữa hai li độ. Tìm các đại lượng}
	\begin{ex} %[3P1B8-1][Loai1]  
		Vật dao động điều hoà. Thời gian ngắn nhất vật đi từ vị trí cân bằng đến li độ $x = 0,5A$ là $0,1\ s$. Chu kì dao động của vật là
		\choice
		{$0,24\ s$}
		{$0,8\ s$}
		{$0,12\ s$}
		{\True $1,2\ s$}
		\loigiai{
			Thời gian ngắn nhất vật đi từ vị trí cân bằng đến li độ $x = 0,5A$ là $\dfrac{T}{12}=0,1 \Rightarrow T = 1,2\ s$.
		}
		%\haidong{10}
	\end{ex}
	%%\loai{Cho khoảng thời gian khoảng cách lớn hơn (hoặc nhỏ hơn) b. Tìm các đại lượng khác}
	\begin{ex} %[3P1K8-3][Loai1]  
		Một chất điểm dao động điều hòa với chu kì T. Biết rằng trong một chu kì, quãng thời gian để khoảng cách từ chất điểm tới vị trí cân bằng lớn hơn $2\ cm$ là $\dfrac{T}{3}$. Biên độ dao động của chất điểm là
		\choice
		{$\dfrac{2}{\sqrt{3}}\ cm$}
		{$2\sqrt{2}\ cm$}
		{\True $\dfrac{4}{\sqrt{3}}\ cm$}
		{$4\ cm$}
		\loigiai{
			\immini{\begin{itemize}
					\item Ta có: $\alpha = \dfrac{2\pi}{3}$.
					\item Ta lại có: $|x| > 2\ cm \Rightarrow x>2\ cm/s$ hoặc $x< -2\ cm/s$\\
					$\Rightarrow$ Để khoảng cách từ chất điểm tới vị trí cân bằng lớn hơn 2 cm thì vật đi từ $M_1$ đến $M_2$ và $M_3$ đến $M_4$ như hình vẽ nên
					\begin{align*}
						\widehat{M_1OM_2}=\widehat{M_3OM_4}=\dfrac{\pi}{3}\Rightarrow 2 = A cos\dfrac{\pi}{6}\Rightarrow A =\dfrac{4}{\sqrt{3}}\ cm.
					\end{align*}
			\end{itemize}}{\begin{tikzpicture}[scale=1,thick,>=stealth']
					\tkzInit[ymin=-3,ymax=3,xmin=-3,xmax=3]\tkzClip
					\tkzDefPoints{-2.5/0/m, 2.5/0/n, 0/0/o, 0/2.5/p, 0/-2.5/q}
					\tkzDefShiftPoint[o](-150:2){1}
					\tkzDefShiftPoint[o](-30:2){2}
					\tkzDefShiftPoint[o](30:2){3}
					\tkzDefShiftPoint[o](150:2){4}
					\tkzDefPointBy[projection=onto m--n](2) \tkzGetPoint{h}
					\tkzDefPointBy[projection=onto m--n](1) \tkzGetPoint{h'}
					\draw(o)circle(2cm);
					\tkzDrawSegments[dashed](1,4 2,3)
					\tkzDrawSegments(p,q)
					\tkzDrawSegments[->](m,n)
					\tkzDrawSegments[blue,->](o,1 o,2 o,4 o,3)
					\tkzDrawPoints[size=3](o,1,h,4,2,3,h')
					\tkzMarkAngles[size=0.7cm,arrows=->](4,o,1 2,o,3)
					\node at (1)[below left]{$M_4$};
					\node at (2)[below right]{$M_1$};
					\node at (3)[above right]{$M_2$};
					\node at (4)[above left]{$M_3$};
					\node at (2,0)[above right]{$A$};
					\node at (h')[below right]{$-2$};
					\node at (h)[below left]{$2$};
		\end{tikzpicture}}}
		
		%\haidong{10}
	\end{ex}
	%%\loai{Thời gian ngắn nhất vật có độ lớn gia tốc}
	\begin{ex}%[3P1K9-2][Loai1]
		Một vật dao động điều hòa theo phương trình $x = 3\cos(5\pi t - 0,5\pi)\ cm$, t tính bằng giây. Thời điểm đầu tiên kể từ $t = 0$ gia tốc của vật có giá trị cực đại là
		\choice
		{$0,10\ s$}
		{\True $0,30\ s$}
		{$0,40\ s$}
		{$0,20\ s$}
		\loigiai{
			\begin{center}
				\begin{tikzpicture}
					\tkzDefPoints{0/0/o, 5/0/a, 4.325/0/b, 3.525/0/c, 2.5/0/d,  -5/0/h, -4.325/0/g, -3.525/0/f, -2.5/0/e, -6/0/x', 6/0/x}
					\tkzDefShiftPoint[a](-90:0.5cm){a1}
					\tkzDefShiftPoint[a](-90:1.2cm){a2}
					\tkzDefShiftPoint[o](-90:0.5cm){o1}
					\tkzDefShiftPoint[h](-90:1.2cm){h2}
					\tkzDrawSegments[blue,->,very thick](x',x)
					\tkzDrawPoints[size=5,fill=black](o,a,b,c,d,e,f,g,h)
					\tkzLabelPoint[above](o){\tiny $O$}
					\tkzLabelPoint[above](a){\tiny$A$}
					\tkzLabelPoint[above](b){\tiny$\dfrac{A\sqrt{3}}{2}$}
					\tkzLabelPoint[above](c){\tiny$\dfrac{A\sqrt{2}}{2}$}
					\tkzLabelPoint[above](d){\tiny$\dfrac{A}{2}$}
					\tkzLabelPoint[above](h){\tiny$-A$}
					\tkzLabelPoint[above](g){\tiny $-\dfrac{A\sqrt{3}}{2}$}
					\tkzLabelPoint[above](f){\tiny$-\dfrac{A\sqrt{2}}{2}$}
					\tkzLabelPoint[above](e){\tiny$-\dfrac{A}{2}$}
					\tkzLabelPoint[above](x){\tiny$x$}
					\tkzDrawSegments[->,very thick,red](o1,a1 a2,h2)
					\tkzDrawSegments[very thick,red](a1,a2)
				\end{tikzpicture}
			\end{center}
			\begin{itemize}
				\item Vậy thời điểm cần tìm là: $t’ = t + \dfrac{3T}{4}= 0,3$ s.
			\end{itemize}
		}
		%\haidong{15}
	\end{ex}
	%\loai{Quãng đường vật đi được sau thời gian $\Delta t$ kể từ thời điểm ban đầu - không cần tách thời gian}
	\begin{ex}%[3P1KD-1][Loai1]
		Một vật nhỏ dao động điều hòa có biên độ A, chu kì dao động T , ở thời điểm ban đầu  $t_0=0$  vật đang ở vị trí biên. Quãng đường mà vật đi được từ thời điểm ban đầu đến thời điểm $t=\dfrac{T}{3}$ là
		\choice
		{\True $\dfrac{3A}{2}$}
		{$\dfrac{A}{3}$}
		{$\dfrac{2A}{3}$}
		{$A$}
		\loigiai{
			\begin{itemize}
				\item Sau $\dfrac{T}{3}$, nếu vật ở biên dương sẽ về $-\dfrac{A}{2}$, nếu vật ở biên âm sẽ tới $\dfrac{A}{2}$ $\Rightarrow$ Quãng đường đi được là $\dfrac{3A}{2}$.
				\item 
				\item 
			\end{itemize}
		}
		%\haidong{15}
	\end{ex}
	%%\loai{Tốc độ trung bình giữa hai thời điểm có li độ và gia tốc cho trước}
	\begin{ex}%[3P1KD-4][Loai1] 
		Một vật nhỏ dao động điều hòa theo một quỹ đạo thẳng dài $20\ cm$ với tần số $0,5\ Hz$. Từ thời điểm vật qua vị trí có li độ $5\ cm$ theo chiều dương đến khi gia tốc của vật đạt giá trị cực tiểu lần thứ hai, vật có tốc độ trung bình là
		\choice
		{$27,0\ cm/s$}
		{$26,7\ cm/s$}
		{\True $19,29\ cm/s$}
		{$27,3\ cm/s$}
		\loigiai{
			\immini{\begin{itemize}
					\item 
					Ta có: $2A = 20 \Rightarrow  A = 10\ cm,\ T = \dfrac{1}{f} = 2\ s$.
					\item Ta biểu diễn vị trí có li độ 5 cm theo chiều dương bằng điểm $M_0$ trên hình vẽ.
					\item Gia tốc của vật có giá trị cực tiểu tại biên dương (vị trí điểm M' trên đường tròn).
					\item Từ hình vẽ, ta suy ra quãng đường vật đi từ thời điểm vật qua vị trí có li độ 5 cm theo chiều dương đến khi gia tốc của vật đạt giá trị cực tiểu lần thứ hai là $s = 5 + 10.4 = 45\ cm$.
					\item Thời gian để vật qua vị trí có li độ 5 cm theo chiều dương đến khi gia tốc của vật đạt giá trị cực tiểu lần thứ hai là $\dfrac{T}{6} + T = \dfrac{7T}{6} = \dfrac{7}{3}\ s$.
			\end{itemize}}{\begin{tikzpicture}[scale=1,thick,>=stealth']
					\tkzInit[ymin=-3,ymax=3,xmin=-3,xmax=3]\tkzClip
					\tkzDefPoints{-2.5/0/m, 2.5/0/n, 0/0/o, 0/2.5/p, 0/-2.5/q}
					\tkzDefShiftPoint[o](-60:2){0}
					\tkzDefShiftPoint[o](0:2){m'}
					\tkzDefPointBy[projection=onto m--n](0) \tkzGetPoint{h}
					\draw(o)circle(2cm);
					\tkzDrawSegments[dashed](0,h)
					\tkzDrawSegments(p,q)
					\tkzDrawSegments[->](m,n)
					\tkzDrawSegments[blue,->](o,0 o,m')
					\tkzDrawPoints[size=3](o,h,0,m')
					\tkzMarkAngles[size=0.7cm,arrows=->](0,o,m')
					\node at (0)[below right]{$M_0$};
					\node at (m')[below right]{$M'$};
					\node at (2,0)[above right]{$10$};
					\node at (h)[above]{$5$};
			\end{tikzpicture}}
			\begin{itemize}
				\item Tốc độ trung bình của vật là $v_{tb} = \dfrac{45}{\dfrac{7}{3}} = 19,29$ m/s.
			\end{itemize}
		}
		%\haidong{25}
	\end{ex}
	%%\loai{Các biểu thức $\omega$, T, f}
	\begin{ex} %[3P1Y3-1][Loai1] 
		Một con lắc lò xo gồm một vật nhỏ khối lượng m và lò xo có độ cứng k. Con lắc dao động điều hòa với tần số góc là
		\choice
		{$2\pi \sqrt{\dfrac{m}{k}}$}
		{$2\pi \sqrt{\dfrac{k}{m}}$}
		{$\sqrt{\dfrac{m}{k}}$}
		{\True $\sqrt{\dfrac{k}{m}}$}
		\loigiai{
			Tần số góc của con lắc lò xo: $\omega  = \sqrt{\dfrac{k}{m}}$.
		}
		%\haidong{5}
	\end{ex}
	\begin{ex}%[3P1-19.1K][Loai1] 
		\immini[thm]{
			Một vật dao động điều hòa trên trục Ox có đồ thị biểu diễn sự phụ thuộc của li độ x vào thời gian t như hình vẽ, pha ban đầu của dao động là
			\choice
			{$10\pi t-\dfrac{\pi }{2}$}
			{$10\pi t+\dfrac{\pi }{2}$}
			{$-\dfrac{\pi }{2}$}
			{\True $\dfrac{\pi }{2}$}
		}{
			\pgfplotsset{width=7.5cm,height=4cm,compat=1.4}
			\begin{tikzpicture}[draw=Mapcolor,scale=1,thick,>=stealth']
				\begin{axis}[
					axis lines=middle,
					xmin=0,xmax=4.5,xstep=1,ymin=-2.16,ymax=2.5,ystep=0.1,
					grid=major, %Có các tùy chọn: minor|major|both|none
					x grid style={dashed},
					y grid style={dashed},
					ytick={-2,0,2},
					xtick={1,2,3,4,5,6},
					xticklabels={,{0,2} },
					yticklabels={},
					xticklabel style={black,below},
					xlabel style={above=3pt},
					xlabel=$t(s)$,
					ylabel style={right=3pt},
					ylabel={$x$}
					]
					\addplot[line width=1.2pt,domain=0:4,samples=200,Mapcolor]{2*cos(deg(0.5*pi*x+pi/2))};
				\end{axis}
			\end{tikzpicture}
		}
		\loigiai{
			Ban đầu vật đi qua vị trí cân bằng theo chiều âm $\varphi_0=0,5\pi$.
		}
		%\haidong{5}
	\end{ex}
	
	%%\loai{Từ chu kì, tần số tính k hoặc m}
	\begin{ex} %[3P1B3-2][Loai1] 
		Một con lắc lò xo gồm vật nhỏ và lò xo nhẹ có độ cứng $10\ N/m$, dao động điều hòa với chu kì riêng $1\ s$. Khối lượng của vật là
		\choice
		{$100\ g$}
		{\True $250\ g$}
		{$200\ g$}
		{$150\ g$}
		\loigiai{
			Ta có: $T=2\pi \sqrt{\dfrac{m}{k}}\Rightarrow m=\dfrac{T^2.k}{4{\pi}^2}=\dfrac{1.10}{4.{\pi}^2}=0,25\ kg=250\ g$.
		}
		
		%\haidong{5}
	\end{ex}
	\begin{ex} %[3P1Y5-1][Loai1] 
		Một chất điểm có khối lượng m đang dao động điều hòa. Khi chất điểm có vận tốc v thì động năng của nó là
		\choice
		{\True $\dfrac{m{v}^{2}}{2}$}
		{$mv^{2}$}
		{$vm^{2}$}
		{$\dfrac{v{m}^{2}}{2}$}
		\loigiai{
			Động năng của chất điểm được xác định bằng biểu thức $W_{\text{đ}}=\dfrac{1}{2}mv^{2}$.
		}
	\end{ex}
	%\haidong{5}
	\begin{ex}%[3P1BF-1][Loai1]
		Tại nơi có gia tốc trọng trường g, một con lắc đơn có sợi dây dài $\ell$ đang dao động điều hoà. Tần số góc dao động của con lắc là
		\choice
		{$\sqrt{\dfrac{\ell}{g}}$}
		{$2\pi \sqrt{\dfrac{g}{\ell}}$}
		{$\dfrac{1}{2\pi}\sqrt{\dfrac{g}{\ell}}$}
		{\True $\sqrt{\dfrac{g}{\ell}}$}
		\loigiai{
			Tần số góc dao động của con lắc là:  $\omega=\sqrt{\dfrac{g}{\ell}}$.
		}
		%\haidong{5}
	\end{ex}
	%%\loai{Tìm chiều dài max, min của lò xo}
	
	%%\loai{Tìm chiều dài lò xo ở vị trí cân bằng}
	\begin{ex}%[3P1B6-2][Loai1]%Bài 15 dạng 1
		Một con lắc lò xo treo thẳng đứng gồm lò xo có dài tự nhiên là $30\ cm$ và vật nhỏ có khối lượng $200\ g$, lò xo có độ cứng $50\ N/m$. Lấy $g = 10\ m/s^2$. Chiều dài lò xo khi vật ở vị trí cân bằng là
		\choice
		{$32\ cm$}
		{\True $34\ cm$}
		{$35\ cm$}
		{$33\ cm$}
		\loigiai{
			Ta có: $ \Delta \ell_0=\dfrac{mg}{k}=  4\ cm \Rightarrow \ell cb = \ell_0 + \Delta \ell_0  = 34\ cm$.
		}
		%\haidong{4}
	\end{ex}
	
	\begin{ex} %[3P1-16.1B][Loai1]  
		
		Một con lắc đơn dao động điều hòa. Năng lượng sẽ thay đổi như thế nào nếu độ cao cực đại của vật tính từ vị trí cân bằng tăng $2$ lần
		\choice
		{\True Tăng $2$ lần}
		{Giảm $2$ lần}
		{Tăng $4$ lần}
		{Giảm $4$ lần}
		\loigiai{
			Ta có: $W = W_{t\max} = mgh\Rightarrow$ h tăng 2 lần thì W tăng 2 lần.
		}
		%\haidong{10}
	\end{ex}
	\begin{ex}%[3P1BF-2][Loai1]%Bài 18
		Một con lắc đơn dao động điều hoà với biên độ góc $\alpha_0 = 0,1\ rad$ tại nơi có gia tốc $g = 10\ m/s^2$. Tại thời điểm ban đầu, vật đi qua vị trí có li độ dài $s=8\sqrt{3}\ cm$ với tốc độ $v = 20\ cm/s$. Chiều dài dây treo vật là
		\choice
		{$80\ cm$}
		{$100\ cm$}
		{\True $160\ cm$}
		{$120\ cm$}
		\loigiai{
			\begin{itemize}
				\item Ta có: $s=\ell \alpha \Rightarrow \alpha =\dfrac{s}{\ell}=\dfrac{0,08\sqrt{3}}{\ell}$.
				\item $\alpha _0^2=\alpha^2+\dfrac{v^2}{g\ell}\Rightarrow 0,1^2=\dfrac{\left( 0,08\sqrt{3}\right)^2}{\ell^2}+\dfrac{0,2^2}{10\ell}\Rightarrow \ell =-1,2\ m$ hoặc $\ell =1,6\ m$.	
			\end{itemize}
		}
		%\haidong{15}
	\end{ex}
\subsubsection*{PHẦN 2. Thí sinh trả lời từ câu 1 đến câu 4. Trong mỗi ý a), b), c), d) ở mỗi câu, thí sinh chọn đúng hoặc sai.}
\setcounter{ex}{0}
\begin{ex}
	Một chất điểm dao động điều hòa với phương trình $x=2 \cos (4 \pi t-0,5 \pi)\ cm$.
	\choiceTF[t]
	{Pha dao động của chất điểm là $0,5 \pi\ rad$}
	{Tốc độ góc của chất điểm là $2\ rad/s$}
	{Chu kì dao động của chất điểm là $1 \ s$}
	{\True Chất điểm dao động trên quỹ đạo có chiều dài $4 \ cm$}
	\loigiai{
		\begin{itemchoice}
			\itemch Sai.
			\itemch Sai.
			\itemch Sai.
			\itemch Đúng.
		\end{\itemchoice}
		
	}
	%\haidong{25}
\end{ex}
\begin{ex}
	Một chất điểm dao động điều hoà theo phương trình $x=8 \cos \left(2 \pi t+\dfrac{\pi}{4}\right)$ ($x$ tính bằng $cm$ và $t$ tính bằng giây).
	\choiceTF[t]
	{\True Quỹ đạo chuyển động của chất điểm có chiều dài là $16 \ cm$}
	{Chu kì dao động của chất điểm là $2 \ s$}
	{Trong $3\ s$ đầu tiên từ thời điểm $t=0$, chất điểm đi qua vị trí có li độ $-4 \ cm$ là $8$ lần}
	{\True Trong $3,5 \ s$ đầu tiên từ thời điểm $t=0$, chất điểm đi qua vị trí có li độ $-3 \ cm$ là  $7$ lần}
	\loigiai{
		\begin{itemchoice}
			\itemch Đúng.
			\itemch Sai.
			\itemch Sai.
			\itemch Đúng.
		\end{\itemchoice}
	}
	%\haidong{25}
\end{ex}
\begin{ex}
	Trong dao động điều hoà luôn có sự chuyển hòa qua lại giữa động năng và thế năng.
	\choiceTF[t]
	{Động năng đạt cực đại tại biên}
	{\True Động năng và thế năng biến thiên với chu kì bằng một nửa chu kì dao động}
	{\True Thế năng đạt cực tiểu tại vị trí cân bằng}
	{Thế năng và động năng của vật biến thiên cùng tần số với tần số của li độ}
	\loigiai{
		\begin{itemchoice}
			\itemch Sai.
			\itemch Đúng.
			\itemch Đúng.
			\itemch Sai.
		\end{\itemchoice}
	}
	%\haidong{25}
\end{ex}
\begin{ex}
	Một vật có khối lượng $m=50 \ g$ được gắn vào lò xo có độ cứng $k=20 \ N/m$ được đặt nằm ngang. Bỏ qua mọi ma sát và lực cản. Kích thích cho vật dao động điều hòa bằng cách truyền cho vật vận tốc $v=1 \ m/s$ ở vị trí cân bằng theo chiều âm.
	\choiceTF[t]
	{\True Ở thời điểm $t=0$ vật ở vị trí lò xo có chiều dài tự nhiên}
	{Tần số góc dao động là $\omega=200\ rad/s$}
	{Phương trình dao động của vật là $x=5 \cos (20 t-\dfrac{\pi}{2})\ cm$}
	{\True Khoảng thời gian giữa hai lần liên tiếp vận tốc vật đạt cực đại là $t=0,1\pi\ s$}
	\loigiai{
		\begin{itemchoice}
			\itemch Đúng.
			\itemch Sai. Tần số góc dao động là $\omega=\sqrt{\dfrac{k}{m}}=20\ rad/s$.
			\itemch Sai. Phương trình dao động của vật là $x=5 \cos (20 t+\pi / 2)\ cm$.
			\itemch $d)$ Đúng. Khoảng thời gian giữa hai lần vận tốc vật đạt cực đại có thể là $t=k \dfrac{T}{2}=k . \dfrac{\pi}{10} khi k=10 \Rightarrow t=\pi\ s$.
		\end{\itemchoice}
	}
	%\haidong{25}
\end{ex}
\subsubsection*{PHẦN 3. Thí sinh trả lời từ câu 1 đến câu 6.}
\setcounter{ex}{0}
\begin{ex}
	Một vật dao động điều hoà với chu kì $0,4\ s$ với chiều dài quỹ đạo là $9\ cm$. Lấy $\pi^2=10$. Gia tốc cực đại bằng bao nhiêu $cm/s^2$?
	\shortans[oly]{1125}
	\loigiai{
		$1125\ cm/s^2$
	}
	%\haidong{30}
\end{ex}
\begin{ex}
	Một vật dao động điều hòa với chu kì $T$ với tốc độ cực đại là $v_{\text {max}}$. Thời gian ngắn nhất vật đi từ điểm mà tốc độ của vật bằng không đến điểm mà tốc độ của vật bằng $\dfrac{\sqrt{2}}{2}v_{\max}$ là bao nhiêu chu kì (làm tròn kết quả đến chữ số hàng phần trăm)?
	\shortans[oly]{0,13}
	\loigiai{
		0,13
	}
	%\haidong{30}
\end{ex}
\begin{ex}
	Một chất điểm dao động điều hòa với phương trình li độ $x=4 \cos (\omega t+\dfrac{\pi}{3}) \ (cm)$ (t tính bằng giây). Thời điểm đầu tiên gia tốc đổi chiều lần thứ hai là $\dfrac{7}{16}\ s$. Chu kì dao động của chất điểm bằng bao nhiêu giây?
	\shortans[oly]{0,75}
	\loigiai{
		$0,75 \ s$
	}
\end{ex}
\begin{ex}
	Một vật nhỏ thực hiện dao động điều hòa theo phương trình $x=10 \sin \left(4 \pi t+\dfrac{\pi}{2}\right)\ cm$ với $t$ tính bằng giây. Động năng của vật đó biến thiên với chu kì bằng bao nhiêu giây?\\
	\shortans[oly]{0,25}
	\loigiai{
		0,25
	}
	%\haidong{30}
\end{ex}
\begin{ex}
	Một chất điểm dao động điều hoà dọc theo trục $O x$ với biên độ $5 \ cm$ và chu kì $2 \ s$. Lấy $\pi^2=10$. Tại thời điểm $t_0$, chất điểm có tốc độ $2,5 \pi\ cm/s$ thì độ lớn gia tốc của chất điểm là bao nhiêu $\ cm/s^2$ (làm tròn kết quả đến chữ số hàng phần mười)?
	\shortans[oly]{43,3} \loigiai{
		$43,3\ cm/s^2$
	}
	%\haidong{30}
\end{ex}
\begin{ex}%[3P1KD-3][Loai1] 
	Một chất điểm dao động điều hòa với biên độ $5\ cm$ và tần số $2\ Hz$. Tốc độ trung bình giữa hai lần liên tiếp chất điểm đi qua vị trí có li độ bằng $0$ là bao nhiêu $cm/s$?\\
	\shortans[oly]{40}
	\loigiai{
		\begin{itemize}
			\item Ta có: $T = \dfrac{1}{f} = 0,5\ s$.
			\item Giữa hai lần liên tiếp chất điểm đi qua vị trí có li độ bằng 0 thì: 
			\begin{align*}
				t = \dfrac{T}{2} = 0,25\ s,\ S = 2A = 10\ cm \Rightarrow  v_{tb} = \dfrac{S}{t} = \dfrac{10}{0,25} = 40\ cm/s.
			\end{align*}
	\end{itemize}}
	%\haidong{27}
\end{ex}
